% 原文 p.55（後半）


\subsubsectionlink{干渉が観測されない問題}
多くの実験では、ある基底の成分状態の量子的重ね合わせが存在することを示す干渉模様を観測できる（2.2.2節）。しかし、より大きな尺度に進むと、そのような干渉模様は通常、消失する。たとえば電子などの微視的粒子で二重スリット実験を行うと、遠方の検出スクリーン上に干渉縞が現れる（第6章の図6.6）。しかし原子や分子で実験すると、干渉模様は普通、急速に消えてしまう。

% 原文 p.56
\hypertarget{srcpage:56}{}\label{srcpage:56}
従来の教科書の説明では、二重スリット実験で中間的・巨視的な物体の干渉模様が観測できないことを、古典的な光波の干渉との類似により説明するのが普通だった。後者では、干渉模様を観測するためには、回折を起こすスリットの間隔が、光の波長程度でなければならないことが、基礎光学でよく知られている。そして同じ議論を、物質による二重スリット実験という量子的場合に適用する。原子などの粒子のde Broglie波長はきわめて短いので、粒子の波長と同程度の幅と間隔を持つスリットを作るのは、実際には不可能である。この見方では、質量を持つ粒子の干渉模様を観測できないのは、粒子の量子的性質を引き出すのを妨げる、実験装置の不十分な「分解能」に由来する。

しかし、これは事情の一部にすぎない。たとえば、二重スリット実験と精神的には似ているが、微視的なスリットを作る必要という障害を避ける実験設定では、中間的な大きさの分子の空間干渉模様を観測できる（実験の詳細は6.2節）。それでも、回折過程に関係しない実験パラメータ、たとえば回折される分子の周囲の空気密度を変えると、干渉模様が減衰する。したがって、実験の回折段階が、干渉模様として観測できる空間的重ね合わせの生成を明らかに許していても、その模様を観測するのを妨げる、別の因子がある。このことは、「干渉を見る」難しさが、重ね合わせをそもそも生成するという問題だけでは説明できないことを、明瞭に示す。

さらに、さまざまな物理系に関する干渉実験の中には、物質の回折実験という分類にまったく入らないものも多い。たとえば、先に簡単に触れたように、逆向きに流れる電流の重ね合わせを作れる実験がある。これは、電流が二方向の間を往復して振動するという形で、時間的な干渉模様を生むはずである（こうした実験は6.3節で詳しく述べる）。それでも、実験条件が重ね合わせ状態の生成を許しているにもかかわらず、干渉を観測できないことが多く、少なくとも、きわめて精密に調整された実験条件が必要になる。この問題を、古典的な光波と量子的な「確率波」の回折の単純な類似で説明できないのは、明らかである。

ましてvon Neumann測定方式の終状態(2.54)からは、巨視的測定装置を含む重ね合わせが、自然界に遍在するはずだと従う。それなら、なぜ周囲の日常世界で、装置の異なる指針位置の間の干渉を観測しないように見えるのか。von Neumann方式は、重ね合わせが微視的な水準から巨視的な水準へ容易に増幅されるはずだと明瞭に示すのに、なぜ中間的・巨視的領域で、何らかの干渉効果を観測するのはそれほど難しいのか。

% 原文 p.57
\hypertarget{srcpage:57}{}\label{srcpage:57}
本質的には、これが干渉の非観測性の問題である。間もなく、デコヒーレンスが、干渉模様の観測される減衰、あるいはそもそもその模様を実験的にまったく観測できないことを、環境自由度との相互作用の結果として説明することにより、この問題への非常に洗練された一般的な答えを与えるのを見る。

\subsubsectionlink{結果の問題}
経験によれば、あらゆる測定は、測定量の確定した値をもたらす。実際、「結果」「値」「量」などの語の定義自体が、この確定性に本質的に依拠する。他方、von Neumann方式で得られる終状態(2.54)は、系と装置の状態の重ね合わせを表す。2.2.1節の議論から、そのような重ね合わせは、系と装置の組合せが実際には成分状態$\lvert s_i\rangle\lvert a_i\rangle$の一つだけにあるが、どれかを私たちが知らないという、状態の古典的集団とは根本的に異なることが分かっている（上の2.4.2節の分析も参照）。したがって、何らかの収縮機構のような追加の物理過程を与えるか、そのような重ね合わせに適切な解釈を与えない限り、最終的な複合状態が与えられたとき、実際の測定の結果として観測される確定した指針位置を、どう説明するかは明らかでない。

この問題はさらに、異なる二つの側面に分けられる。第一に、測定の完了時に、なぜ装置の指針を、異なる指針位置$\lvert a_i\rangle$の重ね合わせとして知覚しないのかという問いに直面する。そもそも、そのような重ね合わせを観測することが何を意味するとしても、なぜ測定には結果があるように見えるのか。第二に、何が特定の結果を「選ぶ」のかを問える。すなわち、測定を実現する実験の各実行で、なぜ他の可能な状態$\lvert a_{j\ne i}\rangle$ではなく、特定の指針位置$i$、したがって特定の指針状態$\lvert a_i\rangle$を観測するのか。両方を合わせて、結果の問題と呼ぶことにする。

結果の問題は、第1章で述べたSchrödingerの猫のパラドックスの直接の基礎にある。その第一の部分が、最終的な原子と猫の複合状態は、古典的には互いに排他的な状態の重ね合わせで記述されると、量子力学が予言するように見えることに関わるのを、思い出そう。簡単のため、ここではハンマーと毒を明示的には議論に含めない。この部分は単に、原子と猫の間のvon Neumann型の、測定に似た相互作用から生じると理解できる。ここで原子は微視的系$\mathcal S$に対応し、猫は、その生死が原子の状態の指標、すなわち「指針」であるという意味で、巨視的装置を表す。式(2.52)に従えば、「測定」方式は次となる。

% 原文 p.58
\hypertarget{srcpage:58}{}\label{srcpage:58}
\begin{align}
\lvert\text{「原子は未崩壊」}\rangle\lvert c_r\rangle&\longrightarrow\lvert\text{「原子は未崩壊」}\rangle\lvert\text{「猫は生存」}\rangle,\tag{2.60a}\\
\lvert\text{「原子は崩壊」}\rangle\lvert c_r\rangle&\longrightarrow\lvert\text{「原子は崩壊」}\rangle\lvert\text{「猫は死亡」}\rangle.\tag{2.60b}
\end{align}
ここで$\lvert c_r\rangle$は猫の初期状態である。量子力学によれば、不安定な原子は常に、原子が崩壊した状態と未崩壊の状態の重ね合わせ、
\begin{equation}\lvert\psi\rangle=\alpha\lvert\text{「原子は未崩壊」}\rangle+\beta\lvert\text{「原子は崩壊」}\rangle\tag{2.61}\end{equation}
で記述される。ここで$\alpha$と$\beta$は、$\lvert\alpha\rvert^2+\lvert\beta\rvert^2=1$を満たす時間依存の係数である。一般的なvon Neumann方式(2.54)とまったく同じく、Schrödinger方程式が線形であることによって、原子と猫の最終的な複合系は、次の形のもつれ状態で記述される。
\begin{equation}\begin{split}\lvert\psi\rangle\lvert c_r\rangle\longrightarrow{}&\alpha\lvert\text{「原子は未崩壊」}\rangle\lvert\text{「猫は生存」}\rangle\\&+\beta\lvert\text{「原子は崩壊」}\rangle\lvert\text{「猫は死亡」}\rangle.\end{split}\tag{2.62}\end{equation}
もはや猫には個別の量子状態を帰属させられず、そのため「猫は生きても死んでもいない」と見える。式(2.62)の右辺の終状態が記述する状況は、しばしばこのように解釈される\footnote{実際、このような状況は、単に私たちの経験の一部でないので、言葉で適切に表すのは難しい。}。ただし「生」と「死」の状態の重ね合わせで記述されるのは、猫自体ではないことに注意しよう。むしろ原子と猫の複合状態が、原子と猫の状態の重ね合わせで表されるのである。ここでもハンマーと毒は省略している。

第1章で説明したように、Schrödingerの猫のパラドックスの第二の部分は、少なくとも量子力学の標準解釈（8.1節）によれば、箱を開けて内容を直接観測し、仮定された「波動関数の収縮」を起こすまで、重ね合わせ(2.62)が存続することに関わる。これは、どうして、またなぜ猫の運命を外部観測者の介入（von Neumannの「erster Eingriff」）に委ねられるのかという問いを生む。もちろん、これは本質的に結果の問題そのものである。重ね合わせ状態(2.62)が与えられたとき、猫の確定した状態、すなわち生きているか死んでいるかの一方にある猫を、観測者が経験することを、どう説明できるだろうか。

結果の問題を解消すること、すなわち、重ね合わせ原理の最重要かつ実験で確認された役割と、測定で単一の確定した結果が観測されることとの見かけ上の矛盾を解くことは、量子力学のあらゆる解釈の中心的動機の一つであった。量子力学の標準的、または「正統的」解釈（詳細は8.1節）は、物理量に対応する観測可能量が確定した値を持つのは、系がその観測可能量の固有状態にある場合、かつその場合に限ると規定する。他方、系がその固有状態の重ね合わせで記述される場合には、系の物理状態が、その観測可能量の何らかの確定した値を持つと語ること自体、無意味と見なされる。

% 原文 p.59
\hypertarget{srcpage:59}{}\label{srcpage:59}
この正統的量子力学の規則は、しばしば固有値と固有状態の結び付き、あるいは、より的確に値と固有状態の結び付きと呼ばれる。

しかし、固有値と固有状態の結び付きは、量子力学の構造からも、いかなる経験的制約からも要請されない\cite{ref068}。さらに「厳密な」固有値と固有状態の結び付きという概念は、それ自体の難点を生む。たとえば測定の結果は、典型的には位置空間に局在した指針で記録される。しかし指針は、決して完全には局在しない。すなわち、位置演算子の厳密な固有状態では記述できない。そのような固有状態は、無限の運動量の広がり、したがって系に含まれる無限のエネルギーに対応するため、非物理的だからである。したがって異なる指針位置に対応する状態は、厳密には互いに直交できない。

そのため、観測可能量とその厳密な固有状態の存在に基づき、固有値と固有状態の結び付きを通じて帰属させられる、古典的な「値」という概念は、しばしば弱められるか、完全に放棄されてきた。そのうえで、量子形式と測定概念をある仕方で再解釈するか、量子力学自体に実際の修正を導入する。前者に属する相対状態解釈と様相解釈は、von Neumann方式(2.54)に現れる最終的な系と装置の複合状態を、量子状態が重ね合わせの形を取るにもかかわらず、結果の存在、または少なくともその主観的知覚を説明するように解釈しようとする。後者に属する物理的収縮模型は（2.2.3節ですでに簡単に触れた）、Schrödinger発展のユニタリ性を破り、波動関数を一成分へ「客観的に」還元する、自然界の何らかの根本的機構の存在を仮定する。これらのさまざまな解釈と、デコヒーレンスとの関係は、第8章でさらに詳しく論じる。

一般に、結果の問題は、確率的理論において、何が特定の結果を現実化するのかという問いに根差す。古典的な確率理論では、この問いに答えることは、少なくとも原理上は難しくない。そこでは、確率の概念は、ある問題の扱いを簡単にする、便利な粗視化手順の帰結にすぎない。しかし物理系の根本的な水準では、具体的な結果は、その基礎にある決定論的な物理法則によって完全に指定される。それが、たとえば何十億個もの原子からなる系について、実際にはどれほど複雑に見えるとしても、である。したがって原理上、古典物理学から確率的な側面をすべて取り除ける。

% 原文 p.60
\hypertarget{srcpage:60}{}\label{srcpage:60}
これに対し、2.1節で詳しく論じたように、量子力学は、本質的に確率的な性格を持つように見える。純粋量子状態は、すでに系の物理状態の完全な記述を表し、その発展は決定論的なSchrödinger方程式で与えられる。それでも、各測定時にどの具体的結果が実現するかを決定する根本的機構は存在しない。したがって、結果の問題は量子力学自体にとって根本的である。そのため、この問題の解決にデコヒーレンスが寄与することに期待できる最大の希望は、なぜ測定で、結果の重ね合わせではなく、可能な結果の一つだけが実際に観測されるかを説明することである。しかし、von Neumannの終状態(2.54)では、いかなる可能な結果も形式的にはまったく特別扱いされないのに、なぜ他の結果ではなく特定の結果が観測者に現れるのかという問いは、デコヒーレンスの範囲外にある、量子力学の解釈の根本的問題に属する。第8章で再びこの話題に戻る。
